\section{A quantitative guide to application regimes}\label{sec:scenarios}

\Cref{sec:applications} lists application areas. The dimensionless groups that decide the heat law are now evaluated for six representative scenarios, which shows at a glance which model of this review is relevant. The table is generated by scenario table from the rounded handbook inputs of \cref{tab:materials} and a lag $\tau_q=1$\,ps; $\delta_T=\sqrt{2\chi/\omega}$ is the thermal penetration depth and $\lambda=c_1/f$ the acoustic wavelength.

\begin{table}[ht]\centering\scriptsize
\caption{Dimensionless groups for six scenarios ($T_0=300$\,K, $\tau_q=1$\,ps). The assessment column applies the rule: Fourier is adequate if $\omega\tau_q<0.1$ and $\omega/\omega_*<0.1$; a lag law may matter if $\omega\tau_q<1.5$ but then the kinetic validity of the model must be checked; otherwise the continuum heat law itself is questionable.}\label{tab:scenarios}
\input{tables/scenarios.tex}
\end{table}

\subsection{Reading the table}
\begin{itemize}
\item \emph{Ultrasonic testing and GHz acoustics (aluminum).} With $\omega\tau_q$ between $6\times10^{-5}$ and $6\times10^{-3}$ and $\omega/\omega_*$ between $1.6\times10^{-4}$ and $1.6\times10^{-2}$, classical coupled thermoelasticity suffices. The thermal penetration depth ($1.8\,\mu$m at $10$\,MHz, $0.18\,\mu$m at $1$\,GHz) is far smaller than the acoustic wavelength ($620\,\mu$m and $6.2\,\mu$m), so the wave feels the temperature field as a thin boundary layer at interfaces and as weak bulk damping.
\item \emph{Picosecond ultrasonics (aluminum, 100\,GHz).} Here $\omega/\omega_*=1.6$, $\omega\tau_q=0.63$, and the wavelength ($62$\,nm) is comparable with the thermal penetration depth ($18$\,nm). This is the corner of the regime map (\cref{sec:synthesis}) where lag and coupling effects could be visible, and also where carrier mean free paths rival the wavelength, so any deviation from the classical result needs a kinetic or atomistic cross-check.
\item \emph{Silicon beam, $h=1\,\mu$m.} At the Zener frequency $f_Z=1.4\times10^{8}$\,Hz, $\omega\tau_q\approx10^{-3}$ and the Fourier law with the exact flexural result \cref{eq:LR} is adequate. The thermal penetration depth ($0.45\,\mu$m) is a sizeable fraction of the thickness, which is why the damping peaks here.
\item \emph{Silicon film $20$\,nm and copper film $50$\,nm.} At the thermal relaxation frequency of the film, $\omega\tau_q$ is $2.3$ and $0.46$ and $\omega/\omega_*$ is $2.7$ and $2.6$. These are the regimes in which a continuum heat law with bulk parameters has to be treated with caution: the film thickness is comparable to or below the carrier mean free path, boundary scattering reduces the effective conductivity, and the relaxation frequency of the continuum model is not reliable.
\end{itemize}

\subsection{Consequences for the choice of theory}
\begin{enumerate}
\item The classical coupled theory suffices for essentially every ultrasonic, above-micron MEMS and geophysical application, and it is the baseline in any case.
\item Lag models can matter only in the lower right-hand corner of \cref{tab:scenarios}, at frequencies above about $10^{10}$\,Hz or thicknesses below about $100$\,nm. There the continuum assumptions are also strained, so a claim of a lag effect needs an independent test.
\item The Guyer--Krumhansl and nonlocal models (\cref{sec:kinetic,sec:size}) are the first refinement once size effects are suspected, because they carry a length scale and lag models do not.
\item The electromechanical and magnetic couplings (\cref{sec:multi}) change the stiffness at all frequencies, and are therefore relevant over the whole table, even where the heat law is not.

\begin{summarybox}[Summary]
The regime map is organized by $\omega\tau_q$ (lag relevance) and $\omega/\omega_*$ (coupling relevance). For $\omega\tau_q\ll0.05$ the classical law suffices; for $0.05\lesssim\omega\tau_q\lesssim1$ lagged models matter; for $\omega\tau_q\gtrsim1$ the exact delay/kinetic law is needed. GK and nonlocal models carry a length scale and are the first refinement once size effects are suspected. Electromechanical and magnetic couplings change stiffness at all frequencies. The six-scenario table applies the coarser thresholds stated in its caption.
\end{summarybox}
\end{enumerate}
Boundary integral equation formulation for thermoelasticity with two relaxation times was developed by \citep{Anwar1994}; reciprocal and variational theorems on the two-temperature Green--Lindsay thermoelasticity theory were given by \citep{Sachan2022}.
